\chapter{Worked Examples}

\section{Purpose of Worked Examples}

This chapter gives explicit toy calculations for the invariants introduced in
the previous chapters. Each example follows the same pattern:

\begin{enumerate}
\item identify the state space;
\item identify the entropy or unresolved defect;
\item identify the admissible refinement step;
\item compute or bound the per-step capacity;
\item derive the resulting depth or rigidity estimate;
\item state the boundary of the example.
\end{enumerate}

The examples are deliberately elementary. Their purpose is to show how the
URF bookkeeping works before applying it to more complicated domains.

\section{Finite Search Space Example}

Let \(\Omega_n\) be a finite search space with
\[
|\Omega_n|=2^n.
\]
Using base-two entropy, the initial unresolved entropy is
\[
H_0=\log_2|\Omega_n|=n.
\]

Suppose an admissible refinement step removes at most \(c>0\) bits of
information. Then the EntropyDepth lower bound gives
\[
\ED\ge \frac{H_0}{c}=\frac{n}{c}.
\]

\begin{example}[One-bit refinement]
If \(c=1\), then at least \(n\) admissible refinement steps are required to
isolate one configuration from \(2^n\) possibilities.
\end{example}

\begin{remark}[Boundary]
This is a lower bound for the declared refinement model only. It does not rule
out a different algorithmic model with larger admissible information access.
\end{remark}

\section{Bounded Branching Partition Example}

Let \(P_t\) be a partition of a finite configuration space \(\Omega\), and let
\(B_t\in P_t\) be the block containing the true configuration. Define
\[
H(P_t)=\log |B_t|.
\]

Assume each refinement step splits \(B_t\) into at most \(b\) effectively
distinguishable subblocks. Then the entropy decrease per step is at most
\[
\log b.
\]
If \(|B_0|=N\), then terminal isolation requires
\[
\ED(P_0)\ge \frac{\log N}{\log b}.
\]

\begin{example}[Binary splitting]
If \(b=2\), each step performs at most one binary distinction. If
\(|B_0|=2^n\), then the depth lower bound is \(n\).
\end{example}

\section{Capacity Constraint Example}

A capacity constraint should be read as a one-step bound, not as a universal
realizability condition.

Let \(H(x)\) be the unresolved entropy of a state \(x\), and let \(C\) be the
maximum entropy removed by one admissible step. If
\[
H(x)>C,
\]
then \(x\) need not be impossible. Rather, \(x\) cannot be fully resolved in
one admissible step.

\begin{theorem}[One-Step Capacity Interpretation]
If terminal resolution requires entropy zero and one step removes at most
\(C\) entropy, then a state \(x\) with \(H(x)>C\) requires more than one
admissible step to resolve.
\end{theorem}

\begin{proof}
If \(x\) were resolved in one step, the entropy removed in that step would be
\(H(x)\). This contradicts the assumption that one step removes at most
\(C<H(x)\) entropy.
\end{proof}

\section{Cycle Graph Spectral Example}

Let \(C_n\) be the cycle graph on \(n\) vertices. The graph Laplacian
eigenvalues are
\[
\lambda_k=2-2\cos\left(\frac{2\pi k}{n}\right),
\qquad
k=0,\ldots,n-1.
\]
The first nonzero eigenvalue is
\[
\lambda_1=2-2\cos\left(\frac{2\pi}{n}\right).
\]

For large \(n\), the approximation
\[
1-\cos x\sim \frac{x^2}{2}
\]
gives
\[
\lambda_1\sim \frac{4\pi^2}{n^2}.
\]
Thus the spectral gap of \(C_n\) tends to zero as \(n\to\infty\).

\begin{remark}[Interpretation]
Cycle graphs do not form a uniform spectral expander family. In URF language,
the spectral coercivity available from the Laplacian degenerates with system
size.
\end{remark}

\section{Expander Contrast Example}

Let \(G_n\) be a \(d\)-regular graph family with graph Laplacians \(L_n\).
Suppose there is a uniform constant \(\lambda>0\) such that
\[
\lambda_1(L_n)\ge \lambda
\]
for every \(n\). Then for every \(f\perp\ker L_n\),
\[
\langle f,{L_n}f\rangle\ge \lambda\|f\|^2.
\]

\begin{theorem}[Uniform Coercivity Example]
Under the uniform spectral gap assumption,
\[
\|f\|^2\le \lambda^{-1}\langle f,{L_n}f\rangle
\]
for every \(f\perp\ker L_n\).
\end{theorem}

\begin{proof}
Divide the inequality
\[
\langle f,{L_n}f\rangle\ge \lambda\|f\|^2
\]
by \(\lambda>0\).
\end{proof}

This example shows how a spectral gap becomes a coercive estimate. It does
not, by itself, provide an entropy-depth lower bound. Entropy and capacity
data must still be supplied.

\section{Local Graph Refinement Example}

Let \(G\) be a graph of maximum degree \(\Delta\), and fix a radius \(R\). The
number of vertices in a radius-\(R\) ball is at most
\[
1+\Delta+\Delta^2+\cdots+\Delta^R.
\]
If vertex labels lie in an alphabet of size \(A\), then the number of possible
labeled local views is bounded by
\[
A^{1+\Delta+\cdots+\Delta^R}
\]
up to the additional choice of local adjacency pattern.

Thus, for fixed \(A\), \(\Delta\), and \(R\), the local view has bounded type
complexity independent of the total graph size.

\begin{example}[Bounded local information]
If an admissible refinement rule only sees radius-\(R\) labeled neighborhoods
in a bounded-degree graph, then one local update cannot directly encode an
unbounded amount of global graph information.
\end{example}

\section{Parallel Update Example}

Suppose a parallel round consists of at most \(m\) local updates, and each
local update removes at most \(c_{\mathrm{loc}}\) entropy. Then one parallel
round removes at most
\[
mc_{\mathrm{loc}}
\]
entropy.

If \(m\) is fixed, this gives a uniform round capacity. If \(m\) grows with
the system size, the round capacity also grows, and a constant-capacity lower
bound no longer follows.

\begin{remark}[Accounting convention]
Parallel examples require an explicit accounting convention. One must state
whether depth is counted per local update, per synchronized round, per
processor, or per global refinement operation.
\end{remark}

\section{Finite Detector Example}

Suppose a measurement device has \(B\) distinguishable outcomes. A single
measurement carries at most
\[
\log B
\]
units of outcome information.

If the unresolved operational entropy is \(H_0\), and terminal resolution
requires eliminating that entropy, then the number of admissible measurements
is bounded below by
\[
\frac{H_0}{\log B}.
\]

\begin{remark}[Physical boundary]
This is an information-accounting example only. A physical theorem would also
need dynamics, measurement coupling, noise assumptions, and an admissibility
model.
\end{remark}

\section{Failed Example: Missing Soundness}

Consider a proposed reduction from a source-domain problem to a refinement
model. Suppose the refinement model has entropy \(H_0\) and capacity \(C\), so
that the model has depth lower bound \(H_0/C\). This does not imply a
source-domain lower bound unless every source-domain process being ruled out
is represented by an admissible refinement trajectory.

\begin{example}[Unsound transfer]
If an algorithm solves the source problem by using a global certificate that
is not represented in the refinement model, then the refinement-model lower
bound does not apply to that algorithm.
\end{example}

This is why reduction soundness is a required field in the application
checklist.

\section{Example Audit Table}

\begin{center}
\begin{tabular}{lll}
\toprule
Example & Main invariant & Conclusion \\
\midrule
Finite search & entropy \(n\) & depth at least \(n/c\) \\
Partition refinement & block entropy & depth at least \(\log N/\log b\) \\
Cycle graph & spectral gap & gap tends to zero \\
Expander family & spectral gap & uniform coercivity \\
Local graph update & local type count & bounded local information \\
Finite detector & outcome count & measurement lower bound \\
\bottomrule
\end{tabular}
\end{center}

\section{Boundary}

This chapter provides worked examples of URF bookkeeping. It does not prove
SAT lower bounds, graph-isomorphism lower bounds, physical impossibility
theorems, unrestricted Chronos-RR, unrestricted H4.1/FGL, P vs NP, or any Clay
problem.
